\documentclass{article}
\usepackage{graphicx}
\usepackage{amsmath}
\usepackage{amssymb}
\usepackage{asymptote}

\title{Olympiad Problems}
\author{Mathematics Initiatives in Nepal}
\date{August 16, 2026}

\begin{document}

\maketitle

\section{Questions}

\begin{enumerate}

\item Triangle $MIN$ is isosceles with $MI = MN$. Medians $\overline{IB}$ and
$\overline{NA}$ are perpendicular to each other, and $IB = NA = 12$.
Find, with proof, the area of $\triangle MIN$.

\medskip

\item You are given that $x$ and $y$ are real numbers.

\begin{enumerate}
    \item Given that $x+y$ and $x+y^2$ are rational numbers, is it always
    true that $x$ and $y$ are also rational numbers?

    \item Given that $x+y$, $x+y^2$, and $x+y^3$ are rational numbers,
    is it always true that $x$ and $y$ are also rational numbers?
\end{enumerate}

\medskip

\item For every positive integer $n$, express, in terms of $n$, the number
of $n$-digit positive integers satisfying both of the following conditions:

\begin{itemize}
    \item No two consecutive digits are equal.
    \item The last digit is prime.
\end{itemize}

\medskip

\item Ram and Hari have created a program that,
given a list of $n$ numbers
$a_1, a_2, \dots, a_n$,
calculates a list
$b_1, b_2, \dots, b_n$
such that, for each $1 \le k \le n$, the number $b_k$ is the number of
times $a_k$ appears in the list. For example, the list
\[
1, \,2, \,3, \,1
\]
outputs
\[
2,\,1,\,1,\,2.
\]

Next, Ram asks Hari to run the program on a list of $2083$ numbers. He then
asks him to apply the program to the resulting list, and so on, until a
number greater than or equal to $k$ appears in the list. Find, with proof,
the largest value of $k$ for which, whatever the initial list of $2083$
positive integers is, it is always possible for Hari to do what Ram asked.

\end{enumerate}

\end{document}
